\section{Pure, or a Task}
\label{sec:ch16-pure-or-a-task}
\index{resolvers!pure and non-pure|(}%

While it is building a selection, the compiler decides how each field will be
run. There are four answers, carried on \code{Selection.Strategy}: default,
serial, pure, and batch. The one worth a section is pure, because it is the
reason a number this book has been printing since
chapter~\ref{ch:the-life-of-a-request} does not mean what its label says.

\index{resolvers!the count is a count of tasks}%
Chapter~\ref{ch:strangling-the-monolith} recorded the outline of this when
Reviews went from 146 resolvers to 26 and 120 of the missing ones were still
being resolved. The finding then was that a resolver with no \code{await} in it
becomes a \code{PureFieldDelegate} that runs inline, and that the diagnostic
event the count comes from is raised only inside \code{ResolverTask} and
\code{BatchResolverTask}. That is still exactly true at this tag. The event is
raised at two call sites, both wrapping genuine asynchronous work:
\code{ResolverTask.Execute.cs} and
\code{BatchResolverTask.cs}. A pure field is invoked from a third place, in
\code{ResolverTaskFactory.cs}, which raises nothing:

\begin{minted}[fontsize=\footnotesize,breaklines]{csharp}
if (resolverContext.TryCreatePureContext(
    selection,
    selectionSetType,
    fieldValue,
    parent,
    out var childContext))
{
    // if we have a pure context we can execute out pure resolver.
    resolverResult = selection.PureResolver!(childContext);
    executedSuccessfully = true;
}
\end{minted}

So the timeline line every Mosaic service prints has always counted resolver
tasks rather than fields resolved. Nothing written before
chapter~\ref{ch:strangling-the-monolith} was wrong, because every field on the
path those chapters measured was asynchronous, and the two numbers were the same
number.

\subsection*{Three decisions, not one}

What I got slightly wrong then is the rule. \enquote{No \code{await} in it} is
true by accident of C\# syntax rather than by design, and purity is settled
three separate times, in three different places, by three different things.

\index{resolvers!classified by return type}%
The first is the source generator, at compile time, and it looks at the
\emph{declared return type} and nothing else. A generic task or value task is
classified as a task, unless what is inside it is an async enumerable, which
gets a kind of its own; an \code{IExecutable}, an \code{IQueryable} and a bare
async enumerable each get one too. A method returning void, or a task with no
type inside it, is rejected as invalid. Everything left over falls through to
pure, which is the last line of the method rather than a case it tests for.
Because C\# obliges an \code{async}
method to return one of those, \enquote{no await} and \enquote{not awaitable}
coincide for ordinary code, but the test being applied is on the signature. A
resolver is also disqualified if any of its parameters is not itself pure, which
is a rule I have not chased to the bottom and am therefore not going to describe.

\index{resolvers!middleware defeats purity}%
The second is schema build, in \code{ObjectField.cs}, and it is the one with
teeth:

\begin{minted}[fontsize=\footnotesize,breaklines]{csharp}
bool IsPureContext()
    => skipMiddleware
        || (context.GlobalComponents.Count == 0
            && fieldMiddlewareDefinitions.Count == 0);
\end{minted}

The generated pure delegate is only kept if that returns true. One global field
middleware registered anywhere in the schema and no field in the schema is pure,
because a middleware chain has to run and the inline path is what skips it. The
file states the consequence as an axiom rather than deriving it, in a comment I
enjoyed finding: \code{// by definition, fields with pure resolvers are parallel
executable.}

\index{resolvers!demoted per selection by a directive}%
The third happens per selection, at operation compile time, and reads the query
text rather than the schema. It is in
\code{OperationCompiler.CompileResolver.cs}:

\begin{minted}[fontsize=\footnotesize,breaklines]{csharp}
if (field.PureResolver is not null && selection.Directives.Count == 0)
{
    return field.PureResolver;
}

for (var i = 0; i < selection.Directives.Count; i++)
{
    if (schema.DirectiveTypes.TryGetDirective(selection.Directives[i].Name.Value, out var type)
        && type.Middleware is not null)
    {
        return null;
    }
}

return field.PureResolver;
\end{minted}

Only a directive that has a registered middleware demotes the field, and only
for that one occurrence. Which settles a question worth asking after
section~\ref{sec:ch16-what-compiling-produces}: \code{@skip} and \code{@include}
have no middleware, because they are consumed as bitmask data rather than as
executable steps, so a field carrying \code{@skip} is still resolved inline by a
\code{PureFieldDelegate}. A hand-written logging directive on the same field is
what would push it back onto the full asynchronous path.

\subsection*{Counting it on something small enough to check}

\index{resolvers!measured on the sample}%
Mosaic's numbers are large enough that the distinction is easy to assert and
hard to see. \code{samples/executor-internals} has a type with three plain
fields and one that returns a \code{Task<string>}, and one asynchronous root
field returning three of them:

\begin{minted}[fontsize=\footnotesize,breaklines]{text}
== pure-fields-cost-no-task
   three things, nine pure fields selected -> 1 resolver tasks
   the same plus one async field per thing -> 4 resolver tasks
   PASS
\end{minted}

Nine pure fields cost nothing at all. Adding one awaitable field per thing adds
exactly three, and the one task in the first run is the root field, which is
asynchronous because every case in that sample needs at least one task to exist.
Both numbers are asserted, so a HotChocolate release that changed the
classification would fail the gate rather than quietly making this section
wrong.

\subsection*{Where the mutation rule is enforced}

\index{mutations!serial execution}%
Serial is the strategy with a specification behind it, and it is applied
nowhere near the executor. When the schema is being built, in
\code{ObjectType.Initialization.cs}, every field of the mutation type is
flagged:

\begin{minted}[fontsize=\footnotesize,breaklines]{csharp}
if (((RegisteredType)context).IsMutationType ?? false)
{
    // if this type represents the mutation type we flag all fields as serially executable
    // so that the operation compiler and execution engine will uphold the spec
    // algorithm to execute mutations serially.
    foreach (var field in definition.Fields)
    {
        field.IsParallelExecutable = false;
    }
}
\end{minted}

Because nothing in a schema except the root returns the mutation type, this
marks root mutation fields and nothing else, which is precisely the rule the
specification states and needs no special case anywhere downstream. Serial also
beats pure: a field already marked serial cannot become pure even if it has a
pure resolver, which is stated in the strategy-inference code and is the right
way round.

\medskip
Four strategies, then, decided while the selection is compiled. What acts on
them is a scheduler, and it turns out to be a smaller and stranger object than
its name suggests.

\index{resolvers!pure and non-pure|)}%
